% 'draft' mode can be used to speed up compilation
\documentclass[twoside,final]{hcmut-report}
\usepackage{codespace}

% Encodings
\usepackage{gensymb,textcomp}

% Better tables
\usepackage{array,longtable,multicol,multirow,siunitx,tabularx}

% Better enum
\usepackage{enumitem}

% Graphics
\usepackage{caption,float}

% Add options for figures, like max width, framing, etc.
\usepackage[export]{adjustbox}

% References
\usepackage[nameinlink]{cleveref}

% Math
\usepackage{amsmath,amssymb}

% Tikz for flowchart
\usepackage{tikz}
\usetikzlibrary{shapes.geometric, arrows.meta, positioning}

% Configurations
\coursename{Kiến trúc Máy tính -- CO2008}
\reporttype{Báo cáo Bài tập Lớn}
\title{Lọc Tín Hiệu và Dự Đoán với Bộ Lọc Wiener\\
       \large Filtering and Prediction Signal with Wiener Filter}
\advisor{& [Tên giảng viên hướng dẫn] &}
\stuname{%
  & [Họ và tên] & [MSSV] \\
}

\AtBeginDocument{\counterwithin{lstlisting}{section}}

\begin{document}
\coverpage%

\clearpage

\tableofcontents
\clearpage

% =================================================================
\section{Giới thiệu}
% =================================================================

Bài tập lớn này thuộc môn Kiến trúc Máy tính (CO2008), yêu cầu em
cài đặt bộ lọc Wiener bậc $M = 2$ bằng hợp ngữ MIPS, chạy trên phần mềm
mô phỏng MARS 4.5.

Ý tưởng của bài toán là: cho trước một tín hiệu đầu vào $x(n)$ (thường là
tổng của tín hiệu gốc và nhiễu) và một tín hiệu mong muốn $d(n)$, thiết kế
bộ lọc FIR có 2 hệ số $h_0$, $h_1$ sao cho đầu ra $y(n)$ xấp xỉ $d(n)$
theo tiêu chí sai số bình phương trung bình nhỏ nhất (MMSE). Đây là bài
toán Wiener-Hopf cổ điển trong xử lý tín hiệu số.

Chương trình cần thực hiện các bước sau:
\begin{enumerate}
  \item Đọc dữ liệu từ hai file \texttt{desired.txt} và \texttt{input.txt}.
  \item Kiểm tra kích thước hai tín hiệu có khớp nhau không.
  \item Tính hệ số tối ưu $h_0$, $h_1$ theo phương trình Wiener-Hopf.
  \item Lọc tín hiệu đầu vào để được $y(n)$.
  \item Tính MMSE giữa $d(n)$ và $y(n)$.
  \item Xuất kết quả ra terminal và file \texttt{output.txt} đúng định dạng.
\end{enumerate}

% =================================================================
\section{Cơ sở lý thuyết}
% =================================================================

\subsection{Bộ lọc FIR và bài toán MMSE}

Bộ lọc FIR bậc $M$ có đầu ra:
\begin{equation}
  y(n) = \sum_{k=0}^{M-1} h_k \cdot x(n-k)
\end{equation}

Trong bài này $M = 2$, nên:
\begin{equation}
  y(n) = h_0 \cdot x(n) + h_1 \cdot x(n-1), \quad
  y(0) = h_0 \cdot x(0)
\end{equation}

Tiêu chí tối ưu là tối thiểu hóa sai số bình phương trung bình:
\begin{equation}
  \text{MMSE} = \frac{1}{N}\sum_{n=0}^{N-1}\bigl(d(n) - y(n)\bigr)^2
  \longrightarrow \min
\end{equation}

\subsection{Phương trình Wiener-Hopf}

Lấy đạo hàm điều kiện tối ưu và đặt bằng 0 ta được hệ phương trình
Wiener-Hopf dạng ma trận $\mathbf{R}\,\mathbf{h} = \boldsymbol{\gamma}_d$,
trong đó:

\begin{equation}
  \mathbf{R} = \begin{pmatrix} r_0 & r_1 \\ r_1 & r_0 \end{pmatrix}, \quad
  \boldsymbol{\gamma}_d = \begin{pmatrix} g_{d0} \\ g_{d1} \end{pmatrix}
\end{equation}

Các đại lượng tương quan được tính từ dữ liệu như sau:
\begin{align}
  r_0    &= \frac{1}{N}\sum_{n=0}^{N-1} x(n)^2
           & \text{(tự tương quan bậc 0)} \\[4pt]
  r_1    &= \frac{1}{N}\sum_{n=1}^{N-1} x(n)\,x(n-1)
           & \text{(tự tương quan bậc 1)} \\[4pt]
  g_{d0} &= \frac{1}{N}\sum_{n=0}^{N-1} d(n)\,x(n)
           & \text{(tương quan chéo bậc 0)} \\[4pt]
  g_{d1} &= \frac{1}{N}\sum_{n=1}^{N-1} d(n)\,x(n-1)
           & \text{(tương quan chéo bậc 1)}
\end{align}

Nghiệm tối ưu theo quy tắc Cramer:
\begin{equation}
  \det(\mathbf{R}) = r_0^2 - r_1^2
\end{equation}
\begin{align}
  h_0 &= \frac{r_0\,g_{d0} - r_1\,g_{d1}}{\det(\mathbf{R})} \\
  h_1 &= \frac{r_0\,g_{d1} - r_1\,g_{d0}}{\det(\mathbf{R})}
\end{align}

\subsection{Định dạng xuất kết quả}

Dựa theo hình minh họa trong đề bài (Figure 3):
\begin{itemize}
  \item Dòng 1: \texttt{Filtered output:} kèm theo $N$ giá trị, giá trị đầu
        chiếm 8 ký tự, các giá trị tiếp theo chiếm 10 ký tự, tất cả có
        4 chữ số thập phân (căn phải).
  \item Dòng 2: \texttt{MMSE: } kèm theo giá trị MMSE có 4 chữ số thập phân.
  \item Nếu kích thước hai file không khớp, chỉ in: \texttt{Error: size not match}.
  \item Nội dung in ra terminal và ghi vào \texttt{output.txt} phải giống nhau.
\end{itemize}

% =================================================================
\section{Thiết kế và cài đặt}
% =================================================================

\subsection{Lưu đồ thuật toán tổng quát}

\begin{figure}[H]
  \centering
  \begin{tikzpicture}[
    node distance=0.7cm,
    box/.style={rectangle, draw, rounded corners, minimum width=5.5cm,
                minimum height=0.75cm, align=center, font=\small},
    decision/.style={diamond, draw, aspect=2.5, align=center, font=\small,
                     inner sep=1pt},
    arrow/.style={-Latex}
  ]
    \node[box, fill=gray!20] (start) {Bắt đầu};
    \node[box, below=of start] (read) {Đọc \texttt{desired.txt} và \texttt{input.txt}};
    \node[decision, below=of read] (check) {Kích thước\\bằng nhau?};
    \node[box, right=2cm of check, fill=red!15] (err) {In ``Error: size not match''\\vào terminal và file};
    \node[decision, below=of check] (checkN) {$N = 0$?};
    \node[box, below=of checkN] (wiener) {Tính $r_0, r_1, g_{d0}, g_{d1}$\\rồi giải Wiener-Hopf $\Rightarrow h_0, h_1$};
    \node[box, below=of wiener] (filter) {Lọc tín hiệu: $y(n) = h_0 x(n) + h_1 x(n-1)$};
    \node[box, below=of filter] (mmse) {Tính MMSE $= \frac{1}{N}\sum(d(n)-y(n))^2$};
    \node[box, below=of mmse] (print) {In kết quả ra terminal và \texttt{output.txt}};
    \node[box, below=of print, fill=gray!20] (end) {Kết thúc};

    \draw[arrow] (start) -- (read);
    \draw[arrow] (read)  -- (check);
    \draw[arrow] (check) -- node[right, font=\small]{Không} (err);
    \draw[arrow] (check) -- node[left,  font=\small]{Có} (checkN);
    \draw[arrow] (checkN) -- node[right, font=\small]{Có} (err);
    \draw[arrow] (checkN) -- node[left,  font=\small]{Không} (wiener);
    \draw[arrow] (wiener) -- (filter);
    \draw[arrow] (filter) -- (mmse);
    \draw[arrow] (mmse)   -- (print);
    \draw[arrow] (print)  -- (end);
    \draw[arrow] (err.south) -- ++(0,-0.3) -| (end.east);
  \end{tikzpicture}
  \caption{Lưu đồ thuật toán chính của chương trình}
  \label{fig:flowchart}
\end{figure}

\subsection{Cấu trúc các chương trình con}

Chương trình được chia thành 6 chương trình con để dễ quản lý và debug,
tương ứng với từng bước trong lưu đồ trên.

\begin{table}[H]
  \centering
  \caption{Danh sách các chương trình con}
  \label{tab:subroutines}
  \begin{tabularx}{\linewidth}{lX}
    \toprule
    \textbf{Nhãn} & \textbf{Chức năng} \\
    \midrule
    \texttt{main}            & Điều phối luồng chính: đọc file, kiểm tra kích thước,
                               gọi các hàm con, mở/đóng file output \\
    \texttt{parse\_floats}   & Phân tích chuỗi ASCII sang mảng số thực 32-bit
                               (xử lý dấu, phần nguyên, phần thập phân) \\
    \texttt{compute\_wiener} & Tính 4 đại lượng tương quan, định thức và hệ số
                               $h_0$, $h_1$ theo Wiener-Hopf \\
    \texttt{compute\_output} & Tính tín hiệu đầu ra $y(n)$ dùng $h_0$, $h_1$ \\
    \texttt{compute\_mmse}   & Tính sai số MMSE giữa $d(n)$ và $y(n)$ \\
    \texttt{print\_results}  & Xuất kết quả ra terminal và file đúng định dạng \\
    \texttt{fmt\_float\_w}   & Định dạng số thực thành chuỗi có độ rộng cố định \\
    \bottomrule
  \end{tabularx}
\end{table}

\subsection{Vùng dữ liệu (.data) và căn chỉnh bộ nhớ}

Các biến bắt buộc theo đề bài được khai báo đúng tên trong vùng \texttt{.data}.
Điểm cần lưu ý là trong MARS, chỉ thị \texttt{.space} \textbf{không} tự
động căn chỉnh bộ nhớ 4 byte như \texttt{.float} hay \texttt{.word}. Do
các chuỗi tên file trước mảng float có tổng $33$ byte (lẻ, không chia hết
cho 4), cần thêm \texttt{.align~2} để tránh lỗi ``store address not aligned''.

\begin{lstlisting}[language={}, caption={Khai báo các biến bắt buộc trong .data}]
desired_fname:        .asciiz "desired.txt"
input_fname:          .asciiz "input.txt"
output_fname:         .asciiz "output.txt"

.align 2                          # can chinh 4 byte cho mang float
desired_signal:       .space 40  # 10 x float32 -- ten bat buoc theo de
input_signal:         .space 40  # 10 x float32 -- ten bat buoc theo de
optimize_coefficient: .space 8   # h0, h1       -- ten bat buoc theo de
mmse:                 .float 0.0 #               -- ten bat buoc theo de
output_signal:        .space 40  # 10 x float32 -- ten bat buoc theo de
\end{lstlisting}

\subsection{Phân tích số thực từ file (parse\_floats)}

Vì MARS không có sẵn hàm đọc số thực từ file, em tự viết parser
bằng cách đọc file vào buffer ASCII rồi phân tích từng ký tự:

\begin{enumerate}
  \item Bỏ qua khoảng trắng (space, tab, ký tự xuống dòng LF/CR).
  \item Kiểm tra dấu \texttt{'-'} nếu có.
  \item Phần nguyên: cộng dồn chữ số (mỗi bước nhân tích lũy với 10).
  \item Phần thập phân: dùng biến \texttt{scale} bắt đầu từ $0.1$,
        mỗi chữ số thập phân chia \texttt{scale} thêm cho 10.
  \item Giới hạn tối đa 10 phần tử để không ghi tràn vào vùng nhớ liền kề.
\end{enumerate}

\begin{lstlisting}[language={}, caption={Pseudo-code của parse\_floats}]
scale = 1.0 / 10.0    # = 0.1
for each fractional digit d:
    value += d * scale
    scale /= 10.0

# Bao ve: dung truoc neu da du 10 phan tu
if count >= 10: goto done
store value; count++
\end{lstlisting}

\subsection{Tính hệ số Wiener (compute\_wiener)}

Sử dụng thanh ghi dấu phẩy động \texttt{\$f16}--\texttt{\$f22} để lưu các
tổng tương quan. Toàn bộ phép tính dùng lệnh single-precision:
\texttt{mul.s}, \texttt{add.s}, \texttt{sub.s}, \texttt{div.s}.

\begin{lstlisting}[language={}, caption={Pseudo-code tính Wiener-Hopf}]
r0 = sum(x[n]^2, n=0..N-1) / N
r1 = sum(x[n]*x[n-1], n=1..N-1) / N
gd0 = sum(d[n]*x[n], n=0..N-1) / N
gd1 = sum(d[n]*x[n-1], n=1..N-1) / N

det = r0*r0 - r1*r1
h0 = (r0*gd0 - r1*gd1) / det    -> luu vao optimize_coefficient[0]
h1 = (r0*gd1 - r1*gd0) / det    -> luu vao optimize_coefficient[1]
\end{lstlisting}

\subsection{Định dạng số thực (fmt\_float\_w)}

MARS không có syscall in số thực với độ rộng trường tùy ý. Nhóm em
giải quyết bằng cách chuyển về số nguyên:

\begin{enumerate}
  \item $\texttt{rounded\_int} = \lfloor |val| \times 10000 + 0.5 \rfloor$
        (dùng lệnh \texttt{floor.w.s}).
  \item $\texttt{int\_part} = \texttt{rounded\_int} \div 10000$,\quad
        $\texttt{frac\_part} = \texttt{rounded\_int} \bmod 10000$.
  \item Xây dựng chuỗi \texttt{['-'] chữ\_số\_nguyên '.' 4\_chữ\_số\_lẻ}
        trong buffer \texttt{num\_buf}.
  \item Thêm khoảng trắng bên trái để đủ độ rộng trường yêu cầu.
  \item In đồng thời ra console (syscall 4) và file (syscall 15).
\end{enumerate}

\subsection{Mở file và ghi output.txt}

MARS dùng hệ cờ file riêng, khác với Linux POSIX:

\begin{lstlisting}[language={}, caption={Mở file output.txt trong MARS}]
# DUNG: MARS dung gia tri 1 cho che do ghi
li   $a1, 1      # 1 = write-only, tu dong tao/xoa cu
li   $v0, 13     # syscall open
syscall

# SAI (Linux POSIX, khong hoat dong trong MARS):
# li $a1, 0x241  # O_WRONLY | O_CREAT | O_TRUNC
\end{lstlisting}

% =================================================================
\section{Các lỗi gặp phải và cách khắc phục}
% =================================================================

Trong quá trình phát triển và kiểm thử, em gặp 3 lỗi nghiêm trọng
và đã tìm ra nguyên nhân gốc rễ để sửa.

\subsection{Lỗi 1: Crash do địa chỉ bộ nhớ không căn chỉnh}

\textbf{Triệu chứng:} Chương trình dừng ngay khi chạy với thông báo lỗi
\textit{``store address not aligned on word boundary 0x10010021''}.

\textbf{Nguyên nhân:} Ba chuỗi tên file (\texttt{desired.txt},
\texttt{input.txt}, \texttt{output.txt}) có độ dài lần lượt là 12, 10,
11 byte (tổng 33 byte -- không chia hết cho 4). Chỉ thị \texttt{.space}
sau đó bị đặt tại địa chỉ lẻ (0x...21), khiến mọi lệnh \texttt{s.s}
(store float) đều crash.

\textbf{Cách sửa:} Thêm \texttt{.align~2} trước \texttt{desired\_signal}
để MARS tự động chèn padding đến địa chỉ chia hết cho 4.

\subsection{Lỗi 2: Đọc số thực sai (phần thập phân bị nhân 10)}

\textbf{Triệu chứng:} Ví dụ số \texttt{3.6} được nạp vào mảng là
\texttt{9.0}, số \texttt{-1.0} trở thành \texttt{-1.0} (may mắn đúng
vì phần lẻ là 0).

\textbf{Nguyên nhân:} Đoạn code khởi tạo \texttt{scale}:
\begin{lstlisting}[language={}]
    div.s $f3, $f4, $f3   # TINH 10.0 / 10.0 = 1.0  (SAI!)
\end{lstlisting}
Thứ tự toán hạng bị đảo ngược, cho ra \texttt{scale = 1.0} thay vì \texttt{0.1}.
Vì vậy chữ số thập phân đầu tiên được nhân với 1.0, kết quả là
$3 + 6 \times 1.0 = 9.0$ thay vì $3 + 6 \times 0.1 = 3.6$.

\textbf{Cách sửa:} Thêm hằng số \texttt{CONST\_ONE: .float 1.0}, rồi:
\begin{lstlisting}[language={}]
    l.s   $f3, CONST_ONE   # f3 = 1.0
    l.s   $f4, CONST_TEN   # f4 = 10.0
    div.s $f3, $f3, $f4    # f3 = 0.1 (DUNG)
\end{lstlisting}

\subsection{Lỗi 3: Không tạo được file output.txt}

\textbf{Triệu chứng:} Chạy xong không thấy file \texttt{output.txt} xuất
hiện trong thư mục.

\textbf{Nguyên nhân:} Dùng cờ Linux (\texttt{0x241 = O\_WRONLY|O\_CREAT|O\_TRUNC})
cho syscall mở file, nhưng MARS dùng hệ số khác: \texttt{1} = ghi
(tự tạo/xóa cũ), \texttt{0} = đọc.

\textbf{Cách sửa:} Thay \texttt{li \$a1, 0x241} bằng \texttt{li \$a1, 1}.

% =================================================================
\section{Xử lý các trường hợp biên}
% =================================================================

Ngoài luồng chính, em có bổ sung thêm các kiểm tra sau để chương
trình không bị crash hay cho kết quả sai trong các tình huống đặc biệt:

\begin{table}[H]
  \centering
  \caption{Các trường hợp biên và cách xử lý}
  \label{tab:edge-cases}
  \begin{tabularx}{\linewidth}{lXl}
    \toprule
    \textbf{Trường hợp} & \textbf{Mô tả / Nguyên nhân cần xử lý} & \textbf{Kết quả} \\
    \midrule
    Kích thước không khớp
      & Hai file có số phần tử khác nhau
      & In lỗi, thoát \\
    File rỗng (N=0)
      & Kích thước $0 = 0$ qua được size check,
        nhưng \texttt{compute\_wiener} sẽ chia cho $N=0$ (NaN)
      & In lỗi, thoát \\
    N=1
      & Không có mẫu $x(n-1)$, nên $r_1 = 0$, $g_{d1} = 0$;
        bộ lọc thoái hóa về $h_1 = 0$
      & Tính được \\
    Input $=$ Desired
      & Bộ lọc tái tạo hoàn hảo, $y(n) = d(n)$
      & MMSE $= 0.0$ \\
    Toàn giá trị âm
      & Dấu âm được xử lý đúng trong parser
      & Tính được \\
    File có $> 10$ phần tử
      & Mảng chỉ có 40 byte (10 float), ghi thêm sẽ tràn
      & Chỉ đọc 10 phần tử đầu, dừng an toàn \\
    \bottomrule
  \end{tabularx}
\end{table}

Riêng trường hợp $N = 0$: sau khi kiểm tra kích thước bằng nhau, chương
trình kiểm tra thêm \texttt{beqz \$t0, pr\_size\_err} để chuyển sang
khối xử lý lỗi trước khi bước vào \texttt{compute\_wiener}.

% =================================================================
\section{Kết quả kiểm thử}
% =================================================================

\subsection{Bộ testcase được cung cấp (5 testcases)}

Tín hiệu mong muốn \texttt{desired.txt} dùng chung cho cả 5 testcase:
\begin{center}
  \texttt{0.0 \; 3.6 \; 4.6 \; 2.3 \; -1.0 \; -2.3 \; -0.3 \; 3.5 \; 6.3 \; 6.0}
\end{center}

\begin{table}[H]
  \centering
  \caption{Kết quả chạy 5 testcase chính}
  \label{tab:tc-results}
  \begin{tabularx}{\linewidth}{cXc}
    \toprule
    \textbf{TC} & \textbf{Tín hiệu đầu ra (\texttt{Filtered output})} & \textbf{MMSE} \\
    \midrule
    1 & \texttt{1.3994\ \ 2.3826\ \ 3.5935\ \ 1.1295\ -0.9636}\\
      & \texttt{-3.7382\ -2.5666\ \ 4.3414\ \ 5.3050\ \ 4.1137} & 1.8287 \\
    \midrule
    2 & \texttt{-0.1159\ \ 4.0982\ \ 2.3170\ \ 1.8235\ \ 0.9891}\\
      & \texttt{\ 2.3268\ \ 0.3964\ -0.8631\ \ 0.9514\ \ 4.1844} & 8.2490 \\
    \midrule
    3 & \texttt{\ 0.6155\ \ 3.3284\ \ 5.4969\ \ 3.5317\ \ 0.6312}\\
      & \texttt{-0.1889\ \ 1.3101\ \ 3.2672\ \ 5.4439\ \ 4.4333} & 1.5726 \\
    \midrule
    4 & \texttt{-1.4757\ \ 2.9779\ \ 3.0827\ -0.7328\ \ 0.5656}\\
      & \texttt{-3.5762\ -3.1946\ \ 1.0424\ \ 4.1019\ \ 5.8285} & 3.7424 \\
    \midrule
    5 & \texttt{-0.3512\ -2.5119\ \ 2.9043\ \ 1.6105\ \ 0.7544}\\
      & \texttt{-0.6135\ -1.8629\ \ 3.4794\ \ 3.7304\ \ 5.3390} & 5.6235 \\
    \bottomrule
  \end{tabularx}
\end{table}

Cả 5 testcase đều chạy thành công và xuất ra đúng định dạng yêu cầu
(field width 8 cho giá trị đầu, 10 cho các giá trị sau, 4 chữ số thập
phân). Nội dung file \texttt{output.txt} khớp hoàn toàn với nội dung
in ra terminal.

\subsection{Kiểm thử trường hợp biên}

\begin{table}[H]
  \centering
  \caption{Kết quả kiểm thử edge cases}
  \label{tab:ec-results}
  \begin{tabularx}{\linewidth}{cllX}
    \toprule
    \textbf{EC} & \textbf{Đầu vào} & \textbf{Expected} & \textbf{Kết quả thực tế} \\
    \midrule
    1 & 3 desired, 5 input
      & In lỗi
      & \texttt{Error: size not match} ✓ \\
    2 & 2 file rỗng (N=0)
      & In lỗi
      & \texttt{Error: size not match} ✓ \\
    3 & N=1: desired=\{2.0\}, input=\{1.0\}
      & Tính được
      & \texttt{Filtered output:\ \ 2.0000 / MMSE: 0.0000} ✓ \\
    4 & Input $=$ Desired (10 phần tử 1--10)
      & MMSE=0
      & \texttt{MMSE: 0.0000} ✓ \\
    5 & Toàn giá trị âm (5 phần tử)
      & Tính được
      & \texttt{-1.5121\ -2.3257\ -3.1105\ -3.8953\ -4.6801} / \texttt{MMSE: 0.0988} ✓ \\
    6 & TC1 lại sau khi thêm edge case guards
      & Không đổi
      & Kết quả giống TC1, không có regression ✓ \\
    \bottomrule
  \end{tabularx}
\end{table}

% =================================================================
\section{Kết luận}
% =================================================================

Sau quá trình thực hiện, em đã hoàn thành bài tập với đầy đủ các
yêu cầu của đề bài:

\begin{itemize}
  \item Đọc và phân tích số thực từ file văn bản (tự viết parser vì MARS
        không hỗ trợ đọc float trực tiếp từ file).
  \item Tính hệ số tối ưu $h_0$, $h_1$ theo phương trình Wiener-Hopf
        sử dụng hoàn toàn lệnh dấu phẩy động của MIPS.
  \item Lọc tín hiệu và tính MMSE.
  \item Xuất kết quả đúng định dạng yêu cầu (field width, 4 chữ số thập phân).
  \item Xử lý đầy đủ các trường hợp biên: kích thước không khớp, file
        rỗng, tràn mảng.
  \item Ghi kết quả đồng thời ra terminal và file \texttt{output.txt}.
\end{itemize}

Qua bài tập này, em hiểu sâu hơn về các vấn đề thực tế khi lập trình
hợp ngữ mà lập trình bậc cao thường che giấu: căn chỉnh địa chỉ bộ nhớ,
sự khác biệt giữa các syscall của MARS và Linux, quản lý stack frame khi
gọi hàm lồng nhau, và cách biểu diễn số thực IEEE 754 trong thanh ghi
FPU của MIPS.

\end{document}
